% R1 — LLM vs fusion. DRAFT (2026-10-09). Every number is from results/summary_r1_*.csv
% (python -m src.analyze_r1 test) or summary_r1_mechanism_val.csv; pre-registration:
% PROJECT_STATUS §10d, commit 17d2733 (before any R1 output). Needs: booktabs;
% references.bib (araci2019finbert). TODO P8: sec:textguard is the P6 TextGuard section,
% not written yet.

\subsection{Is the vulnerability in the LLM or in early fusion?}
\label{sec:r1}

In our main experiments the only vulnerable forecaster (Model~5, Qwen2.5-1.5B-Instruct) is
also the only early-fusion model, so fusion type and ``is an LLM'' are confounded. We separate
them by holding the LLM fixed and changing how its reading of the report enters the forecast:
the \emph{same} Qwen2.5-1.5B-Instruct snapshot used as Model~5 becomes a frozen text encoder
(last-layer hidden states, mean-pooled) for the late-fusion heads of Models~3/4. Everything
else is identical to Models~3/4: frozen DLinear/PatchTST backbones, a text head
$\hat y = \mathrm{TS}(x) + 0.1\cdot\mathrm{MLP}(z(e))$ trained in a second stage, train-only
standardisation, learning rate chosen on validation, three seeds. As a second, CPU-friendly
encoder we use FinBERT~\cite{araci2019finbert}, a BERT fine-tuned for financial sentiment, i.e.\
trained to care whether quantities went up or down. This gives four new heads:
M3-Q/M4-Q (Qwen) and M3-F/M4-F (FinBERT). Without text every head reproduces its backbone
exactly (DLinear bit-identical; PatchTST to float32 precision), so the text effect is isolated.
The design, a 108-test Holm family and three readings (R1-a/b/c) were pre-registered before any
embedding was computed.

\paragraph{Mechanism check (validation, before training).}
For each encoder we measure how far a flipped report moves the head's (standardised) input,
relative to swapping in an unrelated report from the same domain. A pre-registered rule calls
an encoder \emph{direction-aware} if this ratio is at least twice BERT's.
Table~\ref{tab:r1-mech} shows that mean-pooled Qwen embeddings are as direction-blind as BERT
($1.02\times$), whereas FinBERT's vector moves three times as far under a flip. No encoder
registers exaggeration. The test split gives the same picture (BERT 0.203, FinBERT 0.653,
Qwen 0.223).

\begin{table}[t]
\centering\small
\caption{Validation mechanism check: distance the fake report moves the head's input, as a
fraction of swapping in a different report (pooled over domains; flipped $n{=}133$,
exaggerated $n{=}131$).}
\label{tab:r1-mech}
\begin{tabular}{lcccc}
\toprule
Encoder & flipped & $\times$ BERT & direction-aware? & exaggerated \\
\midrule
BERT (Models 3/4) & 0.209 & 1.00 & -- & 0.210 \\
FinBERT & \textbf{0.644} & \textbf{3.08} & yes & 0.218 \\
Qwen2.5-1.5B (= Model 5) & 0.213 & 1.02 & no & 0.183 \\
\bottomrule
\end{tabular}
\end{table}

\paragraph{Attack results (test).}
All four heads are moved by text (16/16 screening tests reject). Table~\ref{tab:r1-attack}
reports the flip attack. When Model~5's own LLM is moved from the prompt into a late-fusion
embedding, its directional vulnerability almost vanishes. The Qwen heads move $+0.008\sigma$
toward the lie, against $+0.457\sigma$ for Model~5 on the same windows. The pre-registered
early-vs-late contrast ($F'$) is $+0.45\sigma$ for both heads (Holm $p<10^{-6}$). The FinBERT
heads, whose embeddings do carry direction, are steered significantly more than their BERT
counterparts (encoder contrast $+0.054\sigma$ and $+0.031\sigma$, Holm $p<10^{-6}$), and M3-F
turns 19 forecasts toward the lie against 1 away (Flip Rate 7.6\%, Holm $p=0.003$). This is
still an order of magnitude below Model~5. No head's error rises measurably under attack
(damage CIs include 0). Exaggeration moves no head appreciably ($|\mathrm{DSS}|\le0.005\sigma$
pooled).

\begin{table}[t]
\centering\small
\caption{Flip attack on the test split, pooled over the four domains (251 attacked windows;
Model~5: 246 after NaN pairing). DSS in $\sigma$ units with 95\% stratified bootstrap CIs;
Holm over the pre-registered R1 family ($m{=}108$). M3/M4/M5: P5 predictions, unchanged.}
\label{tab:r1-attack}
\begin{tabular}{llcccc}
\toprule
Model & Encoder & DSS [95\% CI] & Holm $p$ & Flip Rate & toward/away \\
\midrule
M3 & BERT & $-0.005$ [$-0.010$, $+0.002$] & (P5) & 0.4\% & 1/5 \\
M4 & BERT & $+0.003$ [$+0.001$, $+0.006$] & (P5) & 0.4\% & 1/0 \\
M3-Q & Qwen (pooled) & $+0.008$ [$+0.002$, $+0.014$] & 1.00 & 0.4\% & 1/1 \\
M4-Q & Qwen (pooled) & $+0.008$ [$+0.004$, $+0.011$] & $2\cdot10^{-5}$ & 1.2\% & 3/2 \\
M3-F & FinBERT & $+0.049$ [$+0.040$, $+0.059$] & $<10^{-6}$ & 7.6\% & 19/1 \\
M4-F & FinBERT & $+0.034$ [$+0.028$, $+0.041$] & $<10^{-6}$ & 3.2\% & 8/2 \\
\midrule
M5 & Qwen (in prompt) & $+0.457$ [$+0.348$, $+0.572$] & (P5) & 23.6\% & 58/4 \\
\bottomrule
\end{tabular}
\end{table}

\paragraph{Pre-registered reading.}
Our rule calls a head \emph{fooled} if its pooled flip DSS is significantly positive after
Holm. M4-Q meets it, so the pre-registered outcome is \textbf{R1-a}: a late head on the LLM's
representation is steered, and because $F'$ rejects for both Qwen heads, \emph{early fusion
amplifies} the effect. We report this outcome as specified, without amending the rule. The
interpretation needs to be stated precisely, though. The mechanism check shows that mean-pooled
Qwen embeddings are direction-blind (validation ratio $1.02\times$ BERT's). R1 therefore compares
\emph{in-prompt reading} with a \emph{pooled representation} of the same LLM. It does not compare
two ways of exploiting an understanding the representation retains. With the LLM held fixed,
moving the report from the prompt into a pooled embedding removes about 98\% of the manipulation
(Model~5 $+0.457\sigma$ vs the Qwen heads $+0.008\sigma$; $F'$ $+0.45\sigma$ for both heads).
\emph{What makes Model~5 manipulable is reading the report inside the prompt.} The shift that
triggers ``fooled'' for M4-Q ($+0.0075\sigma$, about $1/60$ of Model~5's) is not distinguishable
from its BERT twin (encoder contrast $+0.004\sigma$, Holm $p=0.28$), and its Flip Rate is not
significant (3 toward / 2 away). M3-Q is not fooled.

\paragraph{Late fusion is not immune by design.}
In our main experiments the BERT-based Models~3--4 were immune to the flip attack, and the
mechanism check explains why: the frozen, mean-pooled BERT encoder barely distinguishes a report
from its flipped version. FinBERT is the positive control for that explanation. Its embedding
does encode direction ($3.08\times$ BERT's), and the same late-fusion heads, trained only for
accuracy, become attackable. Both FinBERT heads are fooled and are steered significantly more
than their BERT twins (encoder contrast $+0.054\sigma$ and $+0.031\sigma$, Holm $p<10^{-6}$).
M3-F moves $+0.049\sigma$ toward the lie and turns 19 forecasts toward it against 1 away (Flip
Rate 7.6\%, Holm $p=0.003$). This is about $10\times$ less than Model~5. Each FinBERT head is
steered in two of the three change-flip domains (M3-F: Economy, SocialGood; M4-F: Economy,
Traffic). The only significant push \emph{away} from the lie is M4-F on SocialGood
($-0.038\sigma$). Late-fusion immunity is therefore a property of the encoder, not of the fusion
architecture.

\paragraph{A lesson for pre-registration.}
The R1 plan fixed tests, family and decision rules but no smallest effect size of interest. As a
result, a statistically clear but negligible shift ($+0.0075\sigma$) satisfied ``fooled''. We
report the rule's outcome and the magnitudes side by side. Every subsequent pre-registration in
this work states a smallest effect size of interest alongside each decision rule.

\paragraph{Ablations and defence (outside the family).}
Excluding 2020 leaves every conclusion unchanged (M3-F $+0.045$, M4-F $+0.027$, M4-Q
$+0.009\sigma$; $F'$ $+0.43\sigma$ for both heads). The FinBERT effect is concentrated in the
dense-text domains (Economy, Traffic: M3-F $+0.069$, M4-F $+0.077\sigma$). As pre-registered for
fooled heads, we applied TextGuard~v1 exactly as frozen in Section~\ref{sec:textguard}. It
removes M4-Q's small shift ($+0.008\to-0.001\sigma$) and reduces M4-F's
($+0.034\to+0.025\sigma$), but not M3-F's ($+0.049\to+0.053\sigma$), at the same 14.4\% False
Block Rate as for Model~5. A defence that only works on the LLM it was developed against is not
enough. Any revised gate must also be evaluated on direction-aware late-fusion heads such as
M3-F.

\paragraph{Limitations.}
One pooling (mean over all tokens) and one LLM size are tested. A different readout (last token,
a probe trained on direction) could carry more direction, and larger LLMs are studied in R2. The
late heads are trained for accuracy on small datasets and are scaled by the fusion weight
$w=0.1$, which bounds how far text can move them. Exploratory (not pre-registered): relative to
how far real text moves each model (pooled sensitivity), the FinBERT heads redirect 42\% (M3-F)
and 64\% (M4-F) of that movement toward the lie, against 5--9\% for the Qwen heads and 118\% for
Model~5. The small absolute effects therefore partly reflect the heads' small text budget.
